% Banque d'exercices — chapitre 11
% Le chapitre est déterminé par le nom de ce fichier.

\begin{exo}[Cycle de Stirling et régénérateur]{stirling-regenerateur}
\keywords{machines thermiques, cycle de Stirling, régénérateur, rendement de Carnot, isotherme, isochore}

Le cycle de Stirling idéal est composé de quatre étapes pour $n$ moles de gaz parfait entre une source chaude à $T_c$ et une source froide à $T_f$ ($T_c > T_f$), avec $V_1 < V_2$ :
\begin{enumerate}[label=\arabic*.]
  \item $A\to B$ : détente isotherme à $T_c$ de $V_1$ à $V_2$.
  \item $B\to C$ : refroidissement isochore à $V_2$ de $T_c$ à $T_f$ (chaleur $Q_{BC}$ cédée au régénérateur).
  \item $C\to D$ : compression isotherme à $T_f$ de $V_2$ à $V_1$.
  \item $D\to A$ : réchauffement isochore à $V_1$ de $T_f$ à $T_c$ (chaleur reprise du régénérateur).
\end{enumerate}

\begin{enumerate}
  \item Calculer $Q_{AB}$, $Q_{BC}$, $Q_{CD}$, $Q_{DA}$ et les travaux associés.
  \item Montrer que $Q_{BC} + Q_{DA} = 0$ : le régénérateur restitue exactement la chaleur stockée.
  \item En déduire le rendement $\eta$ du cycle. Comparer au rendement de Carnot.
\end{enumerate}

\begin{indication}
Sur les isothermes, $\Delta U = 0$ donc $Q = -W$. Sur les isochores, $W = 0$ donc $Q = \Delta U = nC_v\Delta T$.
\end{indication}

\begin{solution}
\textbf{1.}
$$Q_{AB} = nRT_c\ln\frac{V_2}{V_1} > 0, \quad W_{AB} = -nRT_c\ln\frac{V_2}{V_1}.$$
$$Q_{BC} = nC_v(T_f - T_c) < 0, \quad W_{BC} = 0.$$
$$Q_{CD} = nRT_f\ln\frac{V_1}{V_2} < 0, \quad W_{CD} = -nRT_f\ln\frac{V_1}{V_2}.$$
$$Q_{DA} = nC_v(T_c - T_f) > 0, \quad W_{DA} = 0.$$

\textbf{2.} $Q_{BC} + Q_{DA} = nC_v(T_f - T_c) + nC_v(T_c - T_f) = 0$. Le régénérateur absorbe la chaleur lors de $B\to C$ et la restitue intégralement lors de $D\to A$.

\textbf{3.} Chaleur totale absorbée de la source chaude : $Q_c = Q_{AB} = nRT_c\ln(V_2/V_1)$.
Travail total produit :
$$W_{total} = nR(T_c - T_f)\ln\frac{V_2}{V_1}.$$
Rendement :
$$\eta = \frac{|W_{total}|}{Q_c} = \frac{nR(T_c-T_f)\ln(V_2/V_1)}{nRT_c\ln(V_2/V_1)} = 1 - \frac{T_f}{T_c} = \eta_{\text{Carnot}}.$$
Grâce au régénérateur, le cycle de Stirling atteint le rendement de Carnot — ce qui en fait l'un des cycles moteurs idéaux les plus élégants.
\end{solution}

\end{exo}

\begin{exo}[Cycle frigorifique de Carnot : coefficient de performance]{cycle-frigorifique-cop}
\keywords{cycle récepteur, réfrigérateur, coefficient de performance, cycle de Carnot, source froide, source chaude}

Une mole de gaz parfait décrit un cycle de Carnot inversé entre une source froide à $T_f=270$ K et une source chaude à $T_c=300$ K. Le cycle comporte :
\begin{itemize}
\item une détente isotherme au contact de la source froide, pendant laquelle le volume double et le gaz reçoit la chaleur $Q_f>0$ ;
\item une compression adiabatique réversible qui porte le gaz à $T_c$ ;
\item une compression isotherme au contact de la source chaude, pendant laquelle le gaz lui cède une chaleur de valeur absolue $Q_c>0$ ;
\item une détente adiabatique réversible qui ramène le gaz à son état initial.
\end{itemize}

\begin{enumerate}
\item Représenter qualitativement le cycle dans le diagramme de Clapeyron et justifier qu'il est parcouru dans le sens antihoraire. Quel est le signe du travail total reçu par le gaz ?
\item Calculer la chaleur $Q_f$ prélevée à la source froide. Montrer, à l'aide des lois de Laplace sur les deux adiabatiques, que le rapport des volumes sur l'isotherme chaude est lui aussi égal à $2$, puis déterminer $Q_c$.
\item En déduire le travail $W$ reçu par cycle et vérifier le bilan énergétique $Q_f+W=Q_c$.
\item Le coefficient de performance frigorifique est défini par $\mathrm{COP}=Q_f/W$. Montrer que
\[
\mathrm{COP}_{\mathrm{Carnot}}=\frac{T_f}{T_c-T_f},
\]
puis calculer sa valeur. Pourquoi peut-elle être supérieure à $1$ sans violer le premier principe ?
\end{enumerate}

\begin{indication}
Sur une détente isotherme réversible d'un gaz parfait, la chaleur reçue vaut $nRT\ln(V_f/V_i)$. Les adiabatiques imposent le même rapport volumétrique aux deux isothermes. Sur un cycle, $\Delta U=0$.
\end{indication}

\begin{solution}
\textbf{Question 1.} Sur l'isotherme froide, le gaz se détend ; sur l'isotherme chaude, il se comprime. La branche de compression est située au-dessus de la branche de détente : le cycle est donc antihoraire. Il est récepteur et vérifie $W>0$ avec la convention des énergies reçues par le gaz.

\textbf{Question 2.} Sur l'isotherme froide,
\[
Q_f=RT_f\ln2.
\]
Notons $V_1\to V_2=2V_1$ la détente froide et $V_3\to V_4$ la compression chaude. Les deux lois de Laplace donnent
\[
T_fV_2^{\gamma-1}=T_cV_3^{\gamma-1},
\qquad
T_cV_4^{\gamma-1}=T_fV_1^{\gamma-1}.
\]
En les divisant, on obtient $V_3/V_4=V_2/V_1=2$. La valeur absolue de la chaleur cédée sur l'isotherme chaude vaut donc
\[
Q_c=RT_c\ln2.
\]

\textbf{Question 3.} Comme $\Delta U_{\mathrm{cycle}}=0$, le gaz reçoit $Q_f$ de la source froide, cède $Q_c$ à la source chaude et reçoit le travail
\[
W=Q_c-Q_f=R(T_c-T_f)\ln2.
\]
On vérifie immédiatement $Q_f+W=Q_c$.

\textbf{Question 4.}
\[
\mathrm{COP}=\frac{RT_f\ln2}{R(T_c-T_f)\ln2}
=\frac{T_f}{T_c-T_f}
=\frac{270}{300-270}=9.
\]
Un joule de travail permet ici de transférer neuf joules depuis la source froide ; dix joules sont alors cédés à la source chaude. Le COP n'est pas un rendement de conversion : l'énergie utile $Q_f$ est déplacée, non créée, et le bilan $Q_f+W=Q_c$ reste satisfait.
\end{solution}
\end{exo}

\begin{exo}[Cycle de Carnot complet d'un gaz parfait]{carnot-cycle-complet}
\keywords{cycle de Carnot, machine thermique, gaz parfait, rendement, isotherme, adiabatique}

De l'air ($\gamma = 1{,}40$) parcourt le cycle de Carnot réversible ABCD entre une source chaude $T_c = 500$ K et une source froide $T_f = 300$ K. On donne $n = 1{,}0$ mol et, sur l'isotherme chaude BC, $V_B = 5{,}0$ L et $V_C = 17$ L. Les transformations AB et CD sont adiabatiques réversibles ; BC et DA sont isothermes réversibles.

\begin{enumerate}
\item Calculer $Q_{BC}$ (chaleur reçue de la source chaude) et $Q_{DA}$ (chaleur cédée à la source froide).
\item En utilisant les adiabatiques, montrer que $V_D/V_A = V_C/V_B$ et calculer $V_A$ sachant que $V_D = 24$ L.
\item Calculer le travail net $W$ par cycle et le rendement $\eta$. Comparer à $\eta_C = 1 - T_f/T_c$.
\item Montrer que $\eta$ ne dépend que de $T_c$ et $T_f$, pas du gaz de travail.
\end{enumerate}

\begin{indication}
Isotherme : $Q = nRT\ln(V_{\text{final}}/V_{\text{initial}})$ pour une expansion. Adiabatique : $TV^{\gamma-1} = \text{cste}$, d'où $V_A/V_D = V_C/V_B$.
\end{indication}

\begin{solution}
\textbf{Question 1.} Expansion isotherme à $T_c$ ($B \to C$) : $Q_{BC} = nRT_c\ln(V_C/V_B) = 8{,}314\times500\times\ln(17/5) \approx +5{,}09$ kJ (chaleur absorbée). Compression isotherme à $T_f$ (de $D$ vers $A$) : $Q_{DA} = nRT_f\ln(V_A/V_D) < 0$ ; on calcule $V_A$ à la question 2.

\textbf{Question 2.} Sur AB : $T_f V_A^{\gamma-1} = T_c V_B^{\gamma-1}$, soit $(V_A/V_B)^{\gamma-1} = T_c/T_f$. Sur CD : $T_c V_C^{\gamma-1} = T_f V_D^{\gamma-1}$, soit $(V_D/V_C)^{\gamma-1} = T_c/T_f$. Les deux membres de droite étant identiques, $V_A/V_B = V_D/V_C$, c'est-à-dire $V_A/V_D = V_B/V_C$ (attention à ne pas inverser ce rapport). Donc
\[
V_A = V_D\,\frac{V_B}{V_C} = 24\times\frac{5}{17} \approx 7{,}06\ \text{L},
\]
une valeur cohérente avec le sens de parcours du cycle : $V_A$ doit être compris entre $V_B = 5$ L (après la compression adiabatique $A\to B$, donc $V_A > V_B$) et $V_D = 24$ L (avant la compression isotherme $D\to A$, donc $V_A < V_D$). Alors $Q_{DA} = nRT_f\ln(V_A/V_D) = 8{,}314\times300\times\ln(5/17) \approx -3{,}05$ kJ.

\textbf{Question 3.} Premier principe sur le cycle : le travail net \emph{fourni} par le gaz vaut $|W| = Q_{BC} - |Q_{DA}| \approx 5{,}09 - 3{,}05 = 2{,}04$ kJ. $\eta = |W|/Q_{BC} \approx 0{,}40 = 40\,\% = 1 - T_f/T_c$. \checkmark

\textbf{Question 4.} Les chaleurs sur les isothermes valent $Q = nR T \ln r$ ; leur rapport $Q_{DA}/Q_{BC} = -T_f/T_c$ indépendamment de $r$, $n$ et $\gamma$. D'où $\eta = 1 + Q_{DA}/Q_{BC} = 1 - T_f/T_c$ : universalité du rendement de Carnot.
\end{solution}
\end{exo}

\begin{exo}[Cycle de Rankine simplifié]{cycle-rankine}
\keywords{Rankine, machine à vapeur, cycle ditherme, rendement, changement de phase}

Un cycle de Rankine idéal pour la vapeur d'eau comprend : (1) compression liquide adiabatique $1\to2$ ; (2) chauffage isobare $2\to3$ jusqu'à la vaporisation complète ; (3) détente adiabatique réversible dans une turbine $3\to4$ ; (4) condensation isobare à $T_f$ $4\to1$.

\begin{enumerate}
\item Pourquoi la compression du liquide (étape $1\to2$) consomme-t-elle peu de travail comparé à la détente vapeur ($3\to4$) ?
\item On donne $h_2 = 420$ kJ·kg$^{-1}$, $h_3 = 3200$ kJ·kg$^{-1}$, $h_4 = 2200$ kJ·kg$^{-1}$, $h_1 = 200$ kJ·kg$^{-1}$ (enthalpies massiques). Calculer le travail net $w = (h_3-h_4) - (h_2-h_1)$ et le rendement $\eta = w/(h_3-h_2)$.
\item Comparer qualitativement $\eta$ au rendement de Carnot entre $T_c \approx 600$ K (vapeur) et $T_f \approx 300$ K.
\end{enumerate}

\begin{indication}
Liquide incompressible : $w_{pompe} = v_l(P_2-P_1)$ très petit. Turbine : $w_T = h_3-h_4$.
\end{indication}

\begin{solution}
\textbf{Question 1.} Le liquide a un volume spécifique $v_l \ll v_g$ : la pompe comprime un liquide quasi-incompressible, $w \approx v_l\Delta P$ est faible. La turbine détend de la vapeur à grand volume : $w_T = h_3-h_4$ est large.

\textbf{Question 2.} $w = (3200-2200) - (420-200) = 1000 - 220 = 780$ kJ·kg$^{-1}$. $\eta = 780/(3200-420) = 780/2780 \approx 28\,\%$.

\textbf{Question 3.} Carnot : $\eta_C = 1 - 300/600 = 50\,\%$. Rankine $\approx 28\,\%$ : inférieur car la détente adiabatique s'écarte de l'isotherme (brouillard diphasique en $4$), et la pompe/condensation dissipent de l'irréversibilité réelle — cycle réaliste des centrales à charbon/nucléaire.
\end{solution}
\end{exo}

\begin{exo}[Équilibre stellaire et durée de vie]{equilibre-stellaire}
\keywords{étoile, équilibre stellaire, Stefan-Boltzmann, pression de radiation, durée de vie, astrophysique}

Une étoile de masse $M$, rayon $R$ et température centrale $T_c$ est modélisée grossièrement ainsi. La pression de radiation vaut $P_{rad} = \tfrac{1}{3}a T^4$ avec $a = 7{,}57\times10^{-16}$ J·m$^{-3}$·K$^{-4}$. La pression thermique du gaz vaut $P_{gas} = n k_B T$. On admet $P_{rad} \approx P_{gas}$ au centre pour une étoile massive.

\begin{enumerate}
\item Écrire la condition d'équilibre hydrostatique $\dfrac{\mathrm{d}P}{\mathrm{d}r} = -\rho \dfrac{GM(r)}{r^2}$ et son ordre de grandeur au centre ($P/R \sim \rho GM/R^2$).
\item Le Soleil ($R = 7\times10^8$ m, $M = 2\times10^{30}$ kg, $T_c \approx 1{,}5\times10^7$ K) rayonne comme un corps noir de surface $T_s = 5800$ K. Calculer sa luminosité $L = 4\pi R^2 \sigma T_s^4$.
\item La réserve d'énergie nucléaire est $\varepsilon \approx 0{,}007\,Mc^2$ (fraction convertie en énergie). Estimer la durée de vie $\tau \approx \varepsilon/L$ (temps si toute cette énergie partait au débit $L$).
\item Pourquoi les étoiles très massives ($M \gg M_\odot$) vivent-elles moins longtemps malgré un réservoir plus grand ?
\end{enumerate}

\begin{indication}
$\sigma = 5{,}67\times10^{-8}$ W·m$^{-2}$·K$^{-4}$. $L$ fortement sensible à $T_s$ ($T^4$). Massives : $L \propto M^{3{,}5}$ environ (loi empirique).
\end{indication}

\begin{solution}
\textbf{Question 1.} L'équilibre hydrostatique oppose la gravité à la poussée thermique : $\mathrm{d}P/\mathrm{d}r \sim -GM\rho/r^2$. Au centre, $\Delta P/R \sim GM\rho/R^2$ : plus la pression centrale est haute, plus l'étoile résiste à la gravité.

\textbf{Question 2.} $L = 4\pi R^2 \sigma T_s^4 = 4\pi (7\times10^8)^2 \times 5{,}67\times10^{-8} \times (5800)^4 \approx 3{,}8\times10^{26}$ W, proche de la constante solaire observée.

\textbf{Question 3.} $\varepsilon \approx 0{,}007 \times 2\times10^{30} \times (3\times10^8)^2 \approx 1{,}3\times10^{45}$ J. Alors
\[
\tau \approx \frac{\varepsilon}{L} \approx \frac{1{,}3\times10^{45}}{3{,}9\times10^{26}} \approx 3{,}2\times10^{18}\ \text{s} \approx 1\times10^{11}\ \text{ans}.
\]
C'est environ dix fois trop long comparé à la durée de vie réellement observée du Soleil ($\sim 10^{10}$ ans) : l'estimation suppose implicitement que \emph{toute} la masse de l'étoile participe à la fusion, alors qu'en réalité seul le cœur -- là où $T$ et $\rho$ sont suffisants pour la fusion, soit grossièrement $10\,\%$ de la masse totale -- brûle de l'hydrogène pendant la séquence principale. En corrigeant par ce facteur, $\tau \approx 0{,}1\,\varepsilon/L \approx 3\times10^{17}$ s $\approx 10^{10}$ ans, cohérent avec l'âge de la séquence principale solaire.

\textbf{Question 4.} $L \propto M^{3{,}5}$ : doubler la masse multiplie la luminosité par $\sim 11$, épuisant le combustible bien plus vite. Les géantes bleues vivent quelques millions d'années seulement.
\end{solution}
\end{exo}


\begin{exo}[Carnot (1824) : rendement universel par l'absurde]{carnot-raisonnement-historique}
\keywords{Carnot, rendement, machine réversible, second principe, calorique, histoire, impossibilité}

Carnot se demande si le rendement d'une machine thermique dépend du fluide de travail ou de la construction mécanique. Il considère deux machines $M_1$ et $M_2$ réversibles fonctionnant entre les mêmes sources chaude ($T_c$) et froide ($T_f$), de rendements respectifs $\eta_1 > \eta_2$. On admet l'\emph{énoncé de Clausius} du second principe : il est impossible qu'un système échange spontanément de la chaleur d'une source froide vers une source chaude sans qu'aucun travail ne soit fourni de l'extérieur.

\begin{enumerate}
\item $M_1$ fonctionne en moteur : elle prélève $Q_c$ à la source chaude et en rejette $Q_f^{(1)} = Q_c(1-\eta_1)$ à la source froide. Quel travail $W_1$ produit-elle par cycle ?
\item On utilise la \emph{totalité} de ce travail $W_1$ pour actionner $M_2$ en sens inverse (cycle récepteur, comme une pompe à chaleur). Exprimer, en fonction de $W_1$ et $\eta_2$, la chaleur $Q_c^{(2)}$ que $M_2$ rejette alors à la source chaude et la chaleur $Q_f^{(2)}$ qu'elle prélève à la source froide.
\item Vérifier que le travail net échangé avec l'extérieur par l'ensemble $\{M_1 + M_2\text{ en sens inverse}\}$ est nul, puis calculer le bilan net de chaleur reçu par la source chaude et par la source froide. Que constate-t-on ?
\item Conclure sur l'existence de $\eta_1 > \eta_2$ pour deux machines réversibles entre mêmes sources. Que dit-on aujourd'hui de ce raisonnement, alors que Carnot croyait au calorique ?
\end{enumerate}

\begin{indication}
Réversibilité : faire tourner $M_2$ en sens inverse échange $Q$ et $W$ avec des signes opposés par rapport à son fonctionnement en moteur, mais les mêmes relations $Q_c^{(2)} = W^{(2)}/\eta_2$ restent valables en valeur absolue. Pour la question 3, faire le bilan sur chaque source séparément et comparer au signe attendu par l'énoncé de Clausius.
\end{indication}

\begin{solution}
\textbf{Question 1.} $W_1 = \eta_1 Q_c$ et $Q_f^{(1)} = Q_c - W_1 = Q_c(1-\eta_1)$.

\textbf{Question 2.} Fonctionnant en sens inverse et consommant un travail $W_1$, $M_2$ vérifie toujours $W_1 = \eta_2\, Q_c^{(2)}$ (même relation qu'en moteur, mais $Q_c^{(2)}$ est maintenant \emph{rejetée} à la source chaude), d'où $Q_c^{(2)} = W_1/\eta_2$. Le premier principe sur son cycle donne $Q_f^{(2)} = Q_c^{(2)} - W_1 = W_1\left(\dfrac{1}{\eta_2} - 1\right) = \dfrac{W_1(1-\eta_2)}{\eta_2}$.

\textbf{Question 3.} L'ensemble reçoit $W_1$ de l'extérieur pour actionner $M_2$, mais $M_1$ restitue ce même $W_1$ à l'extérieur : le bilan de travail net est donc nul, et le dispositif combiné n'échange que de la chaleur avec ses deux sources. Le bilan à la source chaude est
\[
Q_c^{(2)} - Q_c = \frac{W_1}{\eta_2} - Q_c = \frac{\eta_1 Q_c}{\eta_2} - Q_c = Q_c\,\frac{\eta_1-\eta_2}{\eta_2} > 0 :
\]
la source chaude reçoit un flux net de chaleur (elle en reçoit plus de $M_2$ qu'elle n'en cède à $M_1$). Par conservation de l'énergie (travail net nul), la source froide doit céder exactement ce même flux :
\[
Q_f^{(1)} - Q_f^{(2)} = Q_c(1-\eta_1) - \frac{W_1(1-\eta_2)}{\eta_2} = Q_c(1-\eta_1) - \eta_1 Q_c\frac{1-\eta_2}{\eta_2} = -Q_c\,\frac{\eta_1-\eta_2}{\eta_2} < 0.
\]
Sans aucun travail net échangé avec l'extérieur, ce dispositif fait donc passer spontanément une chaleur $Q_c(\eta_1-\eta_2)/\eta_2 > 0$ de la source froide vers la source chaude.

\textbf{Question 4.} C'est très exactement ce que l'énoncé de Clausius interdit. Comme $M_1$ et $M_2$ sont toutes deux réversibles, on peut échanger leurs rôles et refaire le même raisonnement en sens inverse : l'hypothèse $\eta_1 > \eta_2$ (et de même $\eta_2 > \eta_1$) est donc impossible, si bien que $\eta_1 = \eta_2$. Toutes les machines \emph{réversibles} fonctionnant entre les mêmes températures $T_c$ et $T_f$ ont donc nécessairement le même rendement, quels que soient le fluide de travail et les détails de construction. Ce raisonnement par l'absurde, purement logique, reste entièrement valable aujourd'hui : Carnot avait raison sur la structure de l'argument, bien qu'il croyait encore le calorique conservé (une substance matérielle qui se contenterait de « tomber » du chaud vers le froid comme l'eau d'une cascade, sans que la chaleur puisse se convertir en travail). Seule l'expression numérique $\eta_C = 1 - T_f/T_c$ lui manquait, faute d'une définition rigoureuse de la température absolue -- ce que Kelvin fournira vingt ans plus tard (cf. leçon 11).
\end{solution}
\end{exo}


\begin{exo}[Clausius (1850) : réconcilier Carnot et la conservation de l'énergie]{clausius-entropie-argument}
\keywords{Clausius, entropie, second principe, Carnot, conservation de l'énergie, histoire}

Le premier principe autorise a priori un moteur cyclique monotherme convertissant intégralement une chaleur $Q$ en travail ($W = Q$, rendement 100\%). Carnot, lui, affirme qu'aucune machine ne peut avoir un rendement supérieur au cycle réversible entre deux sources $T_c > T_f$. On se propose de retrouver, comme Clausius, une conséquence quantitative de ce résultat.

On modélise le fluide moteur par un gaz parfait de coefficient $\gamma$ parcourant un cycle de Carnot réversible : une détente isotherme à $T_c$ ($A \to B$, volumes $V_A \to V_B$), une détente adiabatique ($B \to C$, la température passant de $T_c$ à $T_f$), une compression isotherme à $T_f$ ($C \to D$, volumes $V_C \to V_D$), puis une compression adiabatique ($D \to A$, ramenant le gaz à $T_c$).

\begin{enumerate}
\item Écrire le bilan énergétique de ce moteur : exprimer le travail $W$ produit par cycle en fonction de la chaleur $Q_c$ reçue à $T_c$ et de la chaleur $Q_f$ rejetée à $T_f$ (comptées positivement), et en déduire $\eta_C = W/Q_c$ en fonction de $Q_c$ et $Q_f$.
\item Exprimer $Q_c$ (reçue sur $A\to B$) et $Q_f$ (rejetée sur $C \to D$) à l'aide de $n$, $R$, $T_c$, $T_f$ et des volumes $V_A, V_B, V_C, V_D$.
\item En appliquant la loi de Laplace $TV^{\gamma-1} = \text{cste}$ sur chacune des deux branches adiabatiques $B \to C$ et $D \to A$, montrer que $V_B/V_A = V_C/V_D$.
\item On définit $q = Q/T$ pour chaque échange réversible (en comptant algébriquement $+Q_c/T_c$ pour la chaleur reçue et $-Q_f/T_f$ pour la chaleur rejetée). En utilisant les questions précédentes, montrer que $\oint \delta q = Q_c/T_c - Q_f/T_f = 0$ sur ce cycle réversible.
\item Un moteur réel de rendement $\eta < \eta_C$ rejette, pour la même chaleur $Q_c$ reçue, une chaleur $Q_f' > Q_f$ au froid (où $Q_f$ est la valeur réversible de la question précédente). Que vaut alors $Q_c/T_c - Q_f'/T_f$ ? Interpréter le signe obtenu à la lumière de l'inégalité de Clausius $\oint \delta Q/T \leq 0$.
\item En quel sens Clausius a-t-il « comblé » la lacune de Carnot tout en conservant son résultat sur le rendement maximal ?
\end{enumerate}

\begin{indication}
Sur une isotherme réversible à température $T$ pour un gaz parfait : $Q = nRT\ln(V_{\text{final}}/V_{\text{initial}})$. Sur une adiabatique réversible, $\delta Q = 0$ et $TV^{\gamma - 1}$ est constant. Pour la question 5, comparer $Q_f'/T_f$ à $Q_f/T_f$ sachant que $Q_c$ est inchangée.
\end{indication}

\begin{solution}
\textbf{Question 1.} Sur un cycle, $\Delta U = 0 = Q_c - Q_f - W$, donc $W = Q_c - Q_f$ et $\eta_C = W/Q_c = 1 - Q_f/Q_c$.

\textbf{Question 2.} Sur l'isotherme $A \to B$ à $T_c$ : $Q_c = nRT_c \ln(V_B/V_A)$ (positif car $V_B > V_A$, le gaz se détend). Sur l'isotherme $C \to D$ à $T_f$, le gaz est comprimé ($V_D < V_C$) et \emph{rejette} la chaleur $Q_f = nRT_f \ln(V_C/V_D) > 0$.

\textbf{Question 3.} Sur $B \to C$ (adiabatique, $T_c \to T_f$) : $T_c V_B^{\gamma-1} = T_f V_C^{\gamma-1}$, soit $(V_B/V_C)^{\gamma - 1} = T_f/T_c$. Sur $D \to A$ (adiabatique, $T_f \to T_c$) : $T_f V_D^{\gamma - 1} = T_c V_A^{\gamma-1}$, soit $(V_A/V_D)^{\gamma-1} = T_f/T_c$. Les deux membres de droite étant égaux, $V_B/V_C = V_A/V_D$, c'est-à-dire $V_B/V_A = V_C/V_D$. Cette égalité ne doit \emph{rien} à l'entropie : elle découle uniquement de la loi de Laplace sur les deux branches adiabatiques.

\textbf{Question 4.} D'après la question 2, $\oint \delta q = \dfrac{Q_c}{T_c} - \dfrac{Q_f}{T_f} = nR\ln(V_B/V_A) - nR\ln(V_C/V_D) = nR\ln\!\left(\dfrac{V_B/V_A}{V_C/V_D}\right)$. Or la question 3 donne $V_B/V_A = V_C/V_D$, donc ce rapport vaut $1$ et $\oint \delta q = nR \ln(1) = 0$.

\textbf{Question 5.} Puisque $Q_c$ est inchangée mais $Q_f' > Q_f$, on a $Q_f'/T_f > Q_f/T_f = Q_c/T_c$ (question 4), donc
\[
\frac{Q_c}{T_c} - \frac{Q_f'}{T_f} < \frac{Q_c}{T_c} - \frac{Q_f}{T_f} = 0.
\]
Pour ce moteur irréversible, $\oint \delta Q/T < 0$ : c'est bien compatible avec (et un cas particulier de) l'inégalité de Clausius $\oint \delta Q/T \leq 0$, l'égalité n'ayant lieu que dans le cas réversible. Le signe négatif traduit le fait qu'une partie de la chaleur prélevée à la source chaude est, à cause de l'irréversibilité, rejetée « en pure perte » à la source froide sans avoir produit de travail supplémentaire : le rendement réel est nécessairement inférieur à $\eta_C$.

\textbf{Question 6.} Clausius accepte la conservation de l'énergie établie par Mayer et Joule, et remarque que la quantité $\oint \delta Q_{\text{rév}}/T$, nulle pour tout cycle réversible (question 4) et strictement négative dès qu'une irréversibilité intervient (question 5), ne dépend que de l'état du système : elle définit une fonction d'état, l'entropie $S$, telle que $dS = \delta Q_{\text{rév}}/T$ le long d'une transformation réversible, et qui ne peut que croître dans un système isolé. Carnot avait donc parfaitement raison sur l'existence d'un rendement maximal universel ; il lui manquait seulement cette grandeur d'état permettant de quantifier, au-delà du seul cas réversible, le degré d'irréversibilité de n'importe quelle transformation.
\end{solution}
\end{exo}
